\documentclass[11pt,a4paper]{article}
\usepackage[utf8]{inputenc}
\usepackage[T1]{fontenc}
\usepackage{amsmath,amssymb,amsthm}
\usepackage{mathtools}
\usepackage{geometry}
\usepackage{titlesec}
\usepackage{fancyhdr}
\usepackage{enumitem}
\usepackage{array}
\usepackage{booktabs}

\geometry{margin=1in}

\theoremstyle{definition}
\newtheorem*{theorem}{Theorem}
\newtheorem*{lemma}{Lemma}

\pagestyle{fancy}
\fancyhf{}
\rhead{Cayley, Collected Papers, Vol.~VII}
\lhead{Pages 351--400}
\cfoot{\thepage}

\title{\textbf{The Collected Mathematical Papers of Arthur Cayley}\\[6pt]
\Large Volume VII\\[12pt]
\large Pages 351--400}
\author{}
\date{}

\begin{document}
\maketitle
\thispagestyle{fancy}

\vspace{-1em}
\noindent\rule{\textwidth}{0.4pt}

%% ========================================================================
%% PAPER 462 (continued): A NINTH MEMOIR ON QUANTICS
%% ========================================================================
\section*{462. A Ninth Memoir on Quantics. (Continued)}
\addcontentsline{toc}{section}{462. A Ninth Memoir on Quantics (pp.~351--353)}

\noindent\textit{[From the Philosophical Transactions of the Royal Society of London, vol.~CLXI (for the year 1871), pp.~17--50. Read November 10, 1870.]}

\medskip
\setcounter{equation}{0}

\noindent\textbf{356.}\quad Reverting to the form
\[
12^\alpha 13^\beta 23^\gamma \ldots f_i f_j \ldots f_k;
\]
if, as before, $n-i$, $n-j$, \&c.\ are each of them $> \frac{1}{2}n$; if there is at least one index $i$ which is $=$ or $< \frac{1}{2}n$ (that is, for which $n-i > \frac{1}{2}n$), and if the order $mn-i-j\ldots$ be $>n$, then the form, or any sum of such forms, is said to be a form or covariant $W$.

Every covariant $W$ is thus a special covariant, but not conversely. In the particular case $m=2$, the form is
\[
12^\alpha f_i f_j \ldots f_k,
\]
which will be a form $W$ if $n-a > \frac{1}{2}n$, or, what is the same thing, $2n-2a > n$, that is if the order be $>n$. Hence, as already mentioned, the covariants $T$ of the degree 2 are $W$, or else $\not{W}$, according as the order is greater than $n$, or as it is equal to or less than $n$.

\smallskip
\noindent\textbf{357. Theorem B.}\quad If any covariant $T$ be expressible as the sum of a form $W$ and of earlier $T$'s than itself, then forming the derivative $(Tf)^2$; either this is not a form $T$, or being a form $T$, it is expressible as the sum of a form $W$ and of earlier $T$'s than itself; or, what is the same thing, $(Tf)^2$, if it be a form $T$, is (like the original $T$) the sum of a form $W$ and of earlier $T$'s than itself. Hence also every form $T$ is the sum of a form $W$, and of forms derived from the functions $X_1$, $X_2$, \ldots\ or, what is the same thing, every covariant whatever is of the form $W + P_X$.

\smallskip
\noindent\textbf{358.}\quad The proof that for a form $f$ of the order $n$ the number of covariants is finite, depends on the assumption that the number is finite for a form $f'$ of the next inferior order $n-1$: this being so, the number of the special covariants of $f$ will be finite; say these are $A_1$, $A_2$, $A_3$, \ldots\ ($f$ is itself one of the series, but we may separate it, and speak of the form $f$ and its special covariants): the forms $W$ are functions of the special covariants, and hence every covariant whatever of $f$ is of the form $F(A) + P_X$; but it requires still a long investigation to pass from this to the theorem of the existence of a finite number of forms $V$ such that every covariant whatever is $F(V)$. I pass this over,

\newpage
\noindent and consider only the covariants of a quantic of the order $n$; these are all of them of the form $W + P_X$.

\smallskip
\noindent\textbf{359.}\quad If in any covariant $W$ the variables be interchanged in any manner whatever, the result is still a covariant $W$; and if from the same or different covariants $W$ we form any linear function with numerical coefficients, the result is still a covariant $W$; but if we form from a covariant $W$ a symbolical product involving any new symbolical coefficients, the result is no longer a covariant $W$, but is (as already mentioned) a sum of a covariant $W$ and of earlier covariants $T$.

\smallskip
\noindent\textbf{360.}\quad It may be remarked that in the case of a single variable, or say of a binary quantic, the only possible index is $i=1$; this is $< \frac{1}{2}n$ except in the case $n=1$; and the covariants $W$ are consequently nonexistent except for the case in question $n=1$; but for $n=1$, they are in fact the only covariants; viz.\ $f=(ax+by)$ is the only covariant, and this is a $W$.

\medskip
\noindent\textbf{The Binary Quantics.}

\smallskip
\noindent\textbf{361.}\quad For the binary quantics the foregoing theorem may be stated as follows:\[6pt]
\textbf{Theorem.}\ \textit{Every covariant of a binary quantic of the order $n$ is expressible as a rational and integral function of the special covariants.} And observing that every rational and integral function of the special covariants may, in virtue of the linear relations which connect them, be expressed as a rational and integral function of a certain finite number of these special covariants, the theorem may be further stated in the form:\[6pt]
\textbf{Theorem.}\ \textit{Every covariant of a binary quantic of the order $n$ is expressible as a rational and integral function of a certain finite number of covariants.}

\smallskip
\noindent\textbf{362.}\quad For the binary quantic of the order $n$, writing as usual
\[
(a, b, \ldots \mathbin{\raisebox{-2pt}{$\overset{\displaystyle\curvearrowright}{\scriptstyle n}$}} \!\! l\,x, y)^n = (ax+by)^n = \ldots,
\]
the number of asympetic covariants is as already mentioned equal to the coefficient of $x^\alpha a^\rho$ in
\[
\frac{1}{1-ax} \cdot \frac{1}{(1-a^2x^2)(1-ax^2)} \cdots \frac{1}{(1-a^nx^n)},
\]
or, what is the same thing, equal to the coefficient of $a^\rho x^\sigma$ in
\[
\frac{1}{1-ax} \cdot \frac{1}{(1-a^2x^2)(1-ax^2)} \cdots \frac{1}{(1-a^nx^n)},
\]
where $\rho$ is the degree and $\sigma$ the order.

\smallskip
\noindent\textbf{363.}\quad For $U$ equal to any one of the asympetic covariants, the degree $\rho$ and the order $\sigma$ satisfy the condition $\rho n - 2\sigma = 0$ or is a positive integer; or, what is the same thing, $\rho n$ and $2\sigma$ are equal or differ by a multiple of $n$, or, again, $\rho$ and $2\sigma$ are equal or differ by a multiple of 2; or, finally, $\rho$ and $\sigma$ are both of them even, or one of them is even and the other odd. In particular, if $\rho$ is even, then $\sigma$ is even; and if $\rho$ is odd, then $\sigma$ is odd. Hence, also, the order of a covariant of an odd degree is odd; and the order of a covariant of an even degree is even.

\smallskip
\noindent\textbf{364.}\quad For the binary quantic of the order $n$, the complete asympetic system of covariants includes the form itself, and it is convenient to speak of the form as a covariant; but in regard to the covariants of a binary quantic, the term is frequently restricted to mean a covariant other than the form itself; and it is in this restricted sense that I shall in general employ the term.

\smallskip
\noindent\textbf{365.}\quad For the several values of $n$, the asympetic covariants (including the form itself) are shown in the following table:\[6pt]
\begin{center}
\begin{tabular}{c|l}
\toprule
\textbf{Order $n$} & \textbf{Asympetic Covariants} \\
\midrule
$n=2$ & $f$ \\
$n=3$ & $f$, $H$, $\Phi$ \\
$n=4$ & $f$, $H$, $i$, $T$ \\
$n=5$ & $f$, $H$, $i$, $T$, $A$, $B$, $C$ \\
$n=6$ & $f$, $H$, $i$, $T$, $A$, $B$, $C$, $D$, $E$ \\
\bottomrule
\end{tabular}
\end{center}

\newpage
\noindent\textbf{366.}\quad For the binary quantic of the order $n$, the number of asympetic covariants is equal to the number of partitions of any number into parts not exceeding $n$; or, what is the same thing, it is equal to the coefficient of $x^\alpha$ in
\[
\frac{1}{(1-x)(1-x^2)\cdots(1-x^n)}.
\]
For $n=2$, this is
\[
\frac{1}{(1-x)(1-x^2)} = 1 + x + 2x^2 + 2x^3 + 3x^4 + \ldots,
\]
where the coefficient of $x^\alpha$ gives the number of asympetic covariants of the order $\alpha$.

\medskip
\noindent\textbf{367.}\quad The generating function for the number of asympetic covariants of the binary quantic of the order $n$ is
\[
\frac{1}{1-ax} \cdot \frac{1}{(1-a^2x^2)(1-ax^2)} \cdots \frac{1}{(1-a^nx^n)},
\]
and the coefficient of $a^\rho x^\sigma$ in the development of this function gives the number of asympetic covariants of the degree $\rho$ and order $\sigma$.

\smallskip
\noindent\textbf{368.}\quad It is to be observed that the asympetic covariants are not in general independent; there exist syzygies, that is, linear relations with numerical coefficients between the products of the asympetic covariants. For example, for the quartic, there is a syzygy connecting $f^3$, $f^2H$, $fH^2$, $H^3$, $i^2$, $iT$, $T^2$; and for the sextic, there are syzygies connecting the products of the asympetic covariants in a more complicated manner.

\smallskip
\noindent\textbf{369.}\quad The question of the independence of the asympetic coviants, or, more generally, the question of the syzygies which exist between the products of the asympetic covariants, is one of great difficulty; and it has been completely solved only for the lower values of $n$. For $n=2$, there are of course no syzygies; for $n=3$, there is one syzygy of degree 12 and order 12; for $n=4$, there are syzygies of degrees 6 and 8; and so on.

\medskip
\noindent\textbf{370.}\quad The theory of the binary quantics may be extended to ternary and higher quantics; but the theory becomes much more complicated, and many of the results which have been established for binary quantics are still wanting for the higher quantics.

\newpage
%% ========================================================================
%% PAPER 463: NOTE ON A DIFFERENTIAL EQUATION
%% ========================================================================
\section*{463. Note on a Differential Equation.}
\addcontentsline{toc}{section}{463. Note on a Differential Equation (pp.~354--356)}

\noindent\textit{[From the Quarterly Journal of Pure and Applied Mathematics, vol.~VIII. (1867), pp.~120--123.]}

\medskip

\noindent\textbf{1.}\quad The differential equation
\[
\frac{d^ny}{dx^n} = x^m y,
\]
where $m$ is a positive integer, has been considered by various authors; and the solution has been expressed in terms of definite integrals, or by means of series. The equation is of considerable importance in Analysis; and it presents itself in various physical problems.

\smallskip
\noindent\textbf{2.}\quad If we assume
\[
y = \int e^{xt} \phi(t) \, dt,
\]
then substituting in the differential equation, we have
\[
\int t^n e^{xt} \phi(t) \, dt = x^m \int e^{xt} \phi(t) \, dt.
\]
Integrating by parts the right-hand side, we obtain
\[
\int t^n e^{xt} \phi(t) \, dt = \int e^{xt} \left(-\frac{d}{dt}\right)^m \phi(t) \, dt;
\]
and this will be satisfied if
\[
t^n \phi(t) = \left(-\frac{d}{dt}\right)^m \phi(t),
\]
or, what is the same thing,
\[
\frac{d^m \phi}{dt^m} = (-1)^m t^n \phi.
\]

\smallskip
\noindent\textbf{3.}\quad The equation for $\phi$ is of the same form as the original equation for $y$, but with the indices $m$ and $n$ interchanged; and the solution of the one equation may be made to depend on the solution of the other. In particular, if $m=1$, the equation for $\phi$ becomes
\[
\frac{d\phi}{dt} = -t^n \phi,
\]
whence
\[
\phi = C e^{-\frac{t^{n+1}}{n+1}}.
\]

\smallskip
\noindent\textbf{4.}\quad Substituting this value of $\phi$, we have
\[
y = C \int e^{xt - \frac{t^{n+1}}{n+1}} \, dt,
\]
where the integral is taken along any path in the plane of the complex variable $t$, which renders the expression finite. The limits of integration may be so chosen that the integrated terms vanish; and the solution thus obtained will be the complete solution, containing $n$ arbitrary constants.

\newpage
\noindent\textbf{5.}\quad For example, if $n=2$, we have
\[
y = C \int e^{xt - \frac{t^3}{3}} \, dt.
\]
The integral may be evaluated by taking the path of integration to be a contour enclosing the origin, or by taking it from $-\infty$ to $\infty$ along the real axis; and the two independent solutions thus obtained will be the two independent integrals of the differential equation $\dfrac{d^2y}{dx^2} = xy$.

\smallskip
\noindent\textbf{6.}\quad The differential equation $\dfrac{d^2y}{dx^2} = xy$ is a particular case of the equation of the parabolic cylinder; and its solutions are connected with the Airy integral and the Bessel functions of order $\pm\frac{1}{3}$. The solutions may be expressed in the form
\[
y = x^{\frac{1}{2}} \left\{ A J_{\frac{1}{3}}\left(\frac{2}{3} x^{\frac{3}{2}}\right) + B J_{-\frac{1}{3}}\left(\frac{2}{3} x^{\frac{3}{2}}\right) \right\},
\]
where $A$ and $B$ are arbitrary constants.

\smallskip
\noindent\textbf{7.}\quad The general equation $\dfrac{d^ny}{dx^n} = x^m y$ may be treated in a similar manner; and the solution may be expressed by means of definite integrals, or by series involving the gamma function. The theory presents many interesting analytical developments; but the complete discussion of the equation would occupy too much space for the present note.

\newpage
%% ========================================================================
%% PAPER 464: NOTE ON PLANA'S LUNAR THEORY
%% ========================================================================
\section*{464. Note on Plana's Lunar Theory.}
\addcontentsline{toc}{section}{464. Note on Plana's Lunar Theory (pp.~357--360)}

\noindent\textit{[From the Monthly Notices of the Royal Astronomical Society, vol.~XXV. (1864--1865), pp.~225--229.]}

\medskip

\noindent The lunar theory of Plana is founded upon a development of the disturbing function in terms of the elements of the lunar orbit; and the expressions for the coordinates of the Moon are obtained by a process of successive approximation. The theory has been carried to a very high degree of accuracy; and the results have been compared with observation with satisfactory agreement.

\smallskip

\noindent In the development of the disturbing function, Plana introduces certain quantities which are functions of the eccentricities and inclinations of the lunar and solar orbits; and these quantities are expanded in series proceeding according to powers of the eccentricities and inclinations. The coefficients of the several terms in these series are expressed by means of the Besselian functions; and the general terms of the series may be written down by means of a known formula.

\smallskip

\noindent The lunar coordinates are expressed in terms of the mean longitude, the mean anomaly, and the mean elongation; and the expressions involve the time explicitly, through the mean motions of the Moon and the Sun. The theory takes account of the secular variations of the lunar elements; and the long-period inequalities are included in the development.

\smallskip

\noindent In the application of the theory to the calculation of the lunar coordinates, it is necessary to substitute numerical values for the elements of the orbit; and the accuracy of the results depends upon the accuracy with which these elements are known. The comparison of the theory with observation affords a means of determining the elements with greater precision; and the corrections thus obtained are applied to the elements in the subsequent calculations.

\newpage
\noindent The disturbing function in Plana's theory is developed in the following manner. Let $r$, $r'$ be the radii vectores of the Moon and Sun respectively; let $v$, $v'$ be their true longitudes; and let $R$ be the distance between them. Then the disturbing function is
\[
\Omega = \frac{m'}{\sqrt{r^2 + r'^2 - 2rr' \cos(v-v')}} - \frac{m'r}{r'^2} \cos(v-v'),
\]
where $m'$ is the mass of the Sun. The first term represents the direct attraction of the Sun upon the Moon; and the second term represents the effect of the Sun's attraction upon the Earth.

\smallskip

\noindent The function $\Omega$ is expanded in a series of cosines of multiples of the argument $v-v'$; and the coefficients of these cosines are functions of the ratio $r/r'$. Since this ratio is small (its mean value being about $\frac{1}{400}$), the coefficients decrease rapidly; and the series is therefore convergent.

\smallskip

\noindent The expansion of $\Omega$ is effected by means of the formula
\[
\frac{1}{\sqrt{1 - 2\alpha \cos\theta + \alpha^2}} = \frac{1}{2} \sum_{n=-\infty}^{\infty} P_n(\cos\theta) \alpha^{|n|},
\]
where $P_n$ is the Legendre's coefficient of order $n$; or, what is the same thing,
\[
\frac{1}{\sqrt{1 - 2\alpha \cos\theta + \alpha^2}} = \sum_{n=0}^{\infty} P_n(\cos\theta) \alpha^n,
\]
for $\alpha < 1$.

\smallskip

\noindent The development of the disturbing function in Plana's theory is one of the most elaborate analytical works which have ever been executed; and it is remarkable for the elegance and completeness of the results. The theory has been verified by comparison with observation; and it forms the foundation of the lunar tables which are employed in the nautical almanacs.

\newpage
%% ========================================================================
%% PAPER 465: NOTE ON THE LUNAR THEORY
%% ========================================================================
\section*{465. Note on the Lunar Theory.}
\addcontentsline{toc}{section}{465. Note on the Lunar Theory (pp.~361--366)}

\noindent\textit{[From the Monthly Notices of the Royal Astronomical Society, vol.~XXV. (1864--1865), pp.~229--237.]}

\medskip

\noindent The lunar theory, as developed by Laplace, Damoiseau, Plana, Hansen, and other investigators, is founded upon the general principles of the Newtonian theory of gravitation; and the method of solution consists in developing the coordinates of the Moon in series proceeding according to powers of the disturbing force. The disturbing force being small (its effect upon the lunar motion being of the order of the ratio of the Sun's parallax to the Moon's parallax, or about $\frac{1}{400}$), the series converge rapidly; and a comparatively small number of terms suffices to give the coordinates with great accuracy.

\smallskip

\noindent The method of the variation of arbitrary constants, which is employed in the lunar theory, consists in assuming for the coordinates expressions of the same form as in the undisturbed elliptic motion, but with the elements of the orbit regarded as variable quantities. These variable elements are determined by differential equations, which express the condition that the coordinates shall satisfy the equations of motion; and the solution of these differential equations gives the variations of the elements produced by the disturbing force.

\smallskip

\noindent The lunar theory presents many interesting and difficult problems; among which may be mentioned the determination of the secular acceleration of the Moon's mean motion, the theory of the long-period inequalities, and the effect of the figure of the Earth upon the lunar motion. These problems have been treated by the great geometers of the last century and the present; and the results have been verified by comparison with the observations of the Moon, which extend over a period of more than two thousand years.

\smallskip

\noindent The following formulae relate to the expressions for the longitude, latitude, and parallax of the Moon in terms of the elements of the orbit. Let $l$ be the mean longitude, $g$ the mean anomaly, $e$ the eccentricity, $\gamma$ the tangent of the inclination, $c$ the mean distance of the node from the mean equinox, and $m$ the ratio of the mean motion of the Sun to that of the Moon. Then the longitude of the Moon may be expressed in the form
\[
v = l + P_1 \sin g + P_2 \sin 2g + P_3 \sin 3g + \ldots,
\]
where $P_1$, $P_2$, $P_3$, \ldots\ are functions of $e$, $\gamma$, $m$, and the other elements.

\newpage
\noindent\textbf{1.}\quad The expressions for the longitude and latitude of the Moon may be written in the following form. Let $\lambda$ be the longitude, $\beta$ the latitude, and $r$ the radius vector; then
\begin{align*}
\lambda &= l + A_1 \sin g + A_2 \sin 2g + \ldots + B_1 \sin(2l-2l') + B_2 \sin(4l-4l') + \ldots \\
&\quad + C_1 \sin(g + 2l - 2l') + C_2 \sin(g - 2l + 2l') + \ldots,\\
\beta &= \gamma \sin(l - \Omega) + D_1 \sin(3l - 2l' - \Omega) + D_2 \sin(l - 2l' + \Omega) + \ldots,
\end{align*}
where $l$ is the mean longitude of the Moon, $l'$ that of the Sun, $g$ the mean anomaly, $\Omega$ the longitude of the node, and $A_1$, $A_2$, \ldots, $B_1$, $B_2$, \ldots\ are coefficients depending on the elements of the orbit.

\smallskip

\noindent\textbf{2.}\quad The coefficients $A_1$, $A_2$, \ldots\ are functions of the eccentricity $e$; and they may be expanded in series proceeding according to powers of $e$. The general term of these series may be expressed by means of the Besselian functions; and the coefficients may be calculated to any required degree of accuracy.

\smallskip

\noindent\textbf{3.}\quad The coefficients $B_1$, $B_2$, \ldots\ relate to the variation and the parallactic inequality; and they are functions of the ratio $m$ of the mean motions. These coefficients have been calculated to a very high degree of accuracy; and the values have been verified by comparison with observation.

\smallskip

\noindent\textbf{4.}\quad The coefficients $C_1$, $C_2$, \ldots\ relate to the evection; and they are functions both of $e$ and of $m$. The evection is the largest of the lunar inequalities; and its theory has been treated with great care by all the investigators of the lunar theory.

\smallskip

\noindent\textbf{5.}\quad The coefficients $D_1$, $D_2$, \ldots\ relate to the reduction to the ecliptic; and they are functions of the inclination and of the ratio $m$. The reduction to the ecliptic is necessary because the latitude of the Moon is measured from the ecliptic, while the elements of the orbit are referred to the mean equator.

\newpage
\noindent\textbf{6.}\quad The expressions for the longitude and latitude involve also terms depending on the solar eccentricity $e'$, and on the action of the planets. These terms are smaller than those depending on the lunar eccentricity; but they are sensible, and they must be taken into account in the construction of accurate lunar tables.

\smallskip

\noindent\textbf{7.}\quad The differential equations which determine the variations of the elements of the lunar orbit are of the form
\begin{align*}
\frac{de}{dt} &= -\frac{a\sqrt{1-e^2}}{\mu e} \frac{\partial R}{\partial \varpi} + \ldots,\\
\frac{d\varpi}{dt} &= \frac{a\sqrt{1-e^2}}{\mu e} \frac{\partial R}{\partial e} + \ldots,
\end{align*}
where $R$ is the disturbing function, $a$ the semi-axis major, $\mu$ the sum of the masses of the Earth and Moon, and $\varpi$ the longitude of the perigee. Similar equations hold for the other elements of the orbit.

\smallskip

\noindent\textbf{8.}\quad The integration of these differential equations gives the secular and periodic variations of the elements. The secular variations are those which increase continually with the time; and the periodic variations are those which oscillate between fixed limits. The secular variations of the eccentricity and inclination are very small; but the secular variation of the mean motion is sensible, and it gives rise to the well-known acceleration of the Moon's mean motion.

\smallskip

\noindent\textbf{9.}\quad The theory of the acceleration of the Moon's mean motion has been the subject of much discussion; and the question whether the observed acceleration can be entirely accounted for by the gravitational action of the Sun has been investigated by many eminent mathematicians. It appears from the most recent researches that the gravitational theory accounts for the greater part of the observed acceleration; but that there remains a small residual which may be due to the tidal friction, or to other causes not yet fully understood.

\newpage
%% ========================================================================
%% PAPER 466: SECOND NOTE ON THE LUNAR THEORY
%% ========================================================================
\section*{466. Second Note on the Lunar Theory.}
\addcontentsline{toc}{section}{466. Second Note on the Lunar Theory (pp.~367--370)}

\noindent\textit{[From the Monthly Notices of the Royal Astronomical Society, vol.~XXV. (1864--1865), pp.~237--245.]}

\medskip

\noindent\textbf{1.}\quad In my former note on the lunar theory, I gave the general expressions for the longitude, latitude, and parallax of the Moon in terms of the elements of the orbit. I now propose to consider more particularly the expressions for the longitude, and to investigate the coefficients of the several inequalities up to the seventh order inclusive.

\smallskip

\noindent\textbf{2.}\quad The longitude of the Moon may be expressed in the form
\[
\lambda = l + \delta\lambda,
\]
where $l$ is the mean longitude, and $\delta\lambda$ is the sum of the several inequalities. The principal inequalities are the following:\[6pt]
\begin{tabular}{ll}
\quad 1.\ The Equation of the centre & $2e \sin g + \frac{5}{4}e^2 \sin 2g + \ldots$ \\
\quad 2.\ The Evection & $-\frac{15}{8}m e \sin(g - 2l + 2l') + \ldots$ \\
\quad 3.\ The Variation & $\frac{15}{4}m \sin(2l - 2l') + \ldots$ \\
\quad 4.\ The Parallactic inequality & $-\frac{15}{8}m^2 \sin(2l - 2l') + \ldots$ \\
\quad 5.\ The Annual equation & $-e' \sin l' + \ldots$ \\
\end{tabular}

\smallskip

\noindent\textbf{3.}\quad The coefficients of these inequalities are functions of the ratio $m$ of the mean motions; and they may be expanded in series proceeding according to powers of $m$. The calculation of these coefficients is one of the most laborious parts of the lunar theory; and it has been carried to a very high degree of accuracy by the various investigators.

\smallskip

\noindent\textbf{4.}\quad The equation of the centre is the inequality depending on the eccentricity of the lunar orbit; and its coefficient is the largest of all the inequalities. The expression for the equation of the centre may be written in the form
\[
\delta\lambda_1 = 2e \sin g + \frac{5}{4}e^2 \sin 2g + \frac{13}{12}e^3 \sin 3g + \ldots
\]
where $g$ is the mean anomaly. The coefficients of this series are known functions of the eccentricity; and they may be calculated to any required degree of accuracy.

\newpage
\noindent\textbf{5.}\quad The evection is the inequality depending on the combined action of the eccentricity and the solar disturbance; and it is the second in magnitude among the lunar inequalities. The expression for the evection may be written in the form
\[
\delta\lambda_2 = -\frac{15}{8}m e \sin(g - 2l + 2l') + \frac{75}{32}m^2 e \sin(g - 4l + 4l') + \ldots
\]
The coefficient of the principal term is $-\frac{15}{8}me$; and the succeeding terms are of higher orders in $m$ and $e$.

\smallskip

\noindent\textbf{6.}\quad The variation is the inequality depending on the solar disturbance alone; and it is the third in magnitude among the lunar inequalities. The expression for the variation may be written in the form
\[
\delta\lambda_3 = \frac{15}{4}m \sin(2l - 2l') + \frac{105}{16}m^2 \sin(4l - 4l') + \ldots
\]
The coefficient of the principal term is $\frac{15}{4}m$; and the succeeding terms are of higher orders in $m$.

\smallskip

\noindent\textbf{7.}\quad The parallactic inequality is the inequality depending on the solar parallax; and it is smaller than the variation. The expression for the parallactic inequality may be written in the form
\[
\delta\lambda_4 = -\frac{15}{8}m^2 \sin(2l - 2l') - \frac{105}{32}m^3 \sin(4l - 4l') + \ldots
\]
The coefficient of the principal term is $-\frac{15}{8}m^2$; and the succeeding terms are of higher orders in $m$.

\smallskip

\noindent\textbf{8.}\quad The annual equation is the inequality depending on the eccentricity of the Earth's orbit; and it is smaller than the parallactic inequality. The expression for the annual equation may be written in the form
\[
\delta\lambda_5 = -e' \sin l' - \frac{1}{2}e'^2 \sin 2l' + \ldots
\]
where $e'$ is the solar eccentricity, and $l'$ is the mean anomaly of the Sun.

\newpage
%% ========================================================================
%% PAPER 467: EXPRESSIONS FOR PLANA'S e, γ IN TERMS OF THE ELLIPTIC e, γ
%% ========================================================================
\section*{467. Expressions for Plana's $e$, $\gamma$ in terms of the Elliptic $e$, $\gamma$.}
\addcontentsline{toc}{section}{467. Expressions for Plana's $e$, $\gamma$ in terms of the Elliptic $e$, $\gamma$ (pp.~371--374)}

\noindent\textit{[From the Monthly Notices of the Royal Astronomical Society, vol.~XXV. (1864--1865), pp.~245--252.]}

\medskip

\noindent\textbf{1.}\quad In Plana's lunar theory, the quantities $e$ and $\gamma$ are not the eccentricity and tangent of the inclination of the actual lunar orbit, but are certain ideal elements, which are connected with the true elements by relations of a somewhat complicated character. The object of the present note is to establish these relations, and to express Plana's $e$ and $\gamma$ in terms of the elliptic (that is, the true) $e$ and $\gamma$.

\smallskip

\noindent\textbf{2.}\quad The relations between Plana's elements and the true elements may be obtained by comparing the expressions for the coordinates in the two theories. In the elliptic theory, the coordinates are expressed in terms of the true eccentricity $e$ and inclination $\gamma$; while in Plana's theory, they are expressed in terms of Plana's $e$ and $\gamma$. By equating the two expressions for the same coordinate, we obtain relations between the two sets of elements.

\smallskip

\noindent\textbf{3.}\quad The comparison gives the following equations:\begin{align*}
e_P &= e - \frac{1}{8}m^2 e + \frac{1}{4}m^2 e' \cos(\varpi - \varpi') + \ldots,\\
\gamma_P &= \gamma - \frac{1}{8}m^2 \gamma + \frac{3}{8}m^2 \gamma' \cos(\Omega - \Omega') + \ldots,
\end{align*}
where $e_P$, $\gamma_P$ are Plana's elements, and $e$, $\gamma$ are the elliptic elements; $e'$, $\gamma'$ are the solar eccentricity and the inclination of the solar orbit; $\varpi$, $\varpi'$ are the longitudes of the perigees; and $\Omega$, $\Omega'$ are the longitudes of the nodes.

\smallskip

\noindent\textbf{4.}\quad These equations show that Plana's $e$ differs from the true eccentricity by quantities of the order $m^2$; and that Plana's $\gamma$ differs from the true inclination by quantities of the same order. The differences are therefore small; but they are sensible, and they must be taken into account in the comparison of the theory with observation.

\newpage
\noindent\textbf{5.}\quad The relations between Plana's elements and the true elements may also be obtained by a direct transformation of the differential equations. If we write the equations of motion in terms of the true elements, and then transform to Plana's elements by means of the above formulae, we obtain the equations of motion in the form given by Plana.

\smallskip

\noindent\textbf{6.}\quad The inverse relations, expressing the true elements in terms of Plana's, are
\begin{align*}
e &= e_P + \frac{1}{8}m^2 e_P - \frac{1}{4}m^2 e' \cos(\varpi - \varpi') + \ldots,\\
\gamma &= \gamma_P + \frac{1}{8}m^2 \gamma_P - \frac{3}{8}m^2 \gamma' \cos(\Omega - \Omega') + \ldots
\end{align*}
These equations are obtained by solving the former equations with respect to $e$ and $\gamma$; and they are of the same degree of accuracy as the former.

\smallskip

\noindent\textbf{7.}\quad The corrections to Plana's elements, depending on the solar eccentricity and the inclination of the solar orbit, are of considerable importance in the lunar theory. They give rise to inequalities in the lunar motion, which are sensible in the comparison of the theory with observation; and they must be taken into account in the construction of accurate lunar tables.

\smallskip

\noindent\textbf{8.}\quad The foregoing results may be extended to the other elements of the lunar orbit; and similar relations may be established between the other elements of Plana's theory and the corresponding true elements. The complete investigation would occupy too much space for the present note; but the method is the same as that which has been employed for the eccentricity and the inclination.

\newpage
%% ========================================================================
%% PAPER 468: ADDITION TO SECOND NOTE ON THE LUNAR THEORY
%% ========================================================================
\section*{468. Addition to Second Note on the Lunar Theory.}
\addcontentsline{toc}{section}{468. Addition to Second Note on the Lunar Theory (pp.~375--376)}

\noindent\textit{[From the Monthly Notices of the Royal Astronomical Society, vol.~XXV. (1864--1865), pp.~252--254.]}

\medskip

\noindent\textbf{1.}\quad In my Second Note on the Lunar Theory, I gave the expressions for the principal inequalities of the Moon's longitude, up to terms of the seventh order inclusive. I now wish to add some remarks on the method of calculating the higher inequalities, and on the degree of accuracy which may be obtained by the theory.

\smallskip

\noindent\textbf{2.}\quad The method of successive approximation, which is employed in the lunar theory, consists in assuming for the coordinates expressions of the form
\[
x = x_0 + x_1 + x_2 + \ldots,
\]
where $x_0$ is the value of the coordinate in the undisturbed elliptic motion, and $x_1$, $x_2$, \ldots\ are the corrections depending on the disturbing force to the first, second, \ldots\ orders. The quantities $x_1$, $x_2$, \ldots\ are determined by differential equations, which are obtained by substituting the assumed expressions in the equations of motion and equating the terms of each order separately.

\smallskip

\noindent\textbf{3.}\quad The calculation of the higher inequalities is attended with considerable labour, on account of the great number of terms which must be taken into account. The number of inequalities of a given order increases rapidly with the order; and the calculation of the coefficients of these inequalities requires the solution of a large number of differential equations. The labour of the calculation has been greatly diminished by the introduction of mechanical methods; but it is still very considerable.

\smallskip

\noindent\textbf{4.}\quad The degree of accuracy which may be obtained by the lunar theory is very high. The coefficients of the principal inequalities are known with great precision; and the errors of the theory are chiefly due to the uncertainty of the elements of the orbit, and to the neglect of the higher inequalities. By carrying the calculation to a sufficient order, the errors of the theory may be made smaller than the errors of observation; and the theory may thus be used to determine the elements of the orbit with the greatest possible accuracy.

\newpage
%% ========================================================================
%% PAPER 469: ON AN EXPRESSION FOR THE ANGULAR DISTANCE OF TWO PLANETS
%% ========================================================================
\section*{469. On an Expression for the Angular Distance of Two Planets.}
\addcontentsline{toc}{section}{469. On an Expression for the Angular Distance of Two Planets (pp.~377--379)}

\noindent\textit{[From the Monthly Notices of the Royal Astronomical Society, vol.~XXV. (1864--1865), pp.~254--258.]}

\medskip

\noindent\textbf{1.}\quad The angular distance of two planets, as seen from the Earth, is an important element in the theory of planetary perturbations; and it is also of use in the determination of the orbits of the planets from observations. The expression for the angular distance may be obtained from the spherical triangle formed by the two planets and the Earth; and it may be developed in a series proceeding according to powers of the ratios of the distances to the radius of the sphere of observation.

\smallskip

\noindent\textbf{2.}\quad Let $r$, $r'$ be the distances of the two planets from the Earth; let $\lambda$, $\lambda'$ be their longitudes; and let $\beta$, $\beta'$ be their latitudes. Then the angular distance $z$ is given by the formula
\[
\cos z = \sin\beta \sin\beta' + \cos\beta \cos\beta' \cos(\lambda - \lambda').
\]
This formula is exact; but it is not convenient for calculation, on account of the number of trigonometrical functions which it involves.

\smallskip

\noindent\textbf{3.}\quad For the purpose of calculation, it is more convenient to have an expression for $\cos z$ in terms of the rectangular coordinates of the planets. If $x$, $y$, $z$ and $x'$, $y'$, $z'$ are the coordinates of the two planets referred to the ecliptic, then
\[
rr' \cos z = xx' + yy' + zz'.
\]
This expression may be developed in a series by substituting for the coordinates their values in terms of the elements of the orbits.

\smallskip

\noindent\textbf{4.}\quad The development of the expression for $\cos z$ gives
\begin{align*}
\cos z &= \cos(\lambda - \lambda') \cos\beta \cos\beta' + \sin\beta \sin\beta' \\
&= \left(1 - \frac{1}{2}\beta^2 + \ldots\right)\left(1 - \frac{1}{2}\beta'^2 + \ldots\right) \cos(\lambda - \lambda') + \beta\beta' + \ldots
\end{align*}
where $\beta$, $\beta'$ are expressed in parts of the radius. This development is convenient when the latitudes are small, as is usually the case with the planets.

\newpage
\noindent\textbf{5.}\quad The expression for the angular distance may also be obtained by means of the theory of spherical coordinates. If $\theta$, $\theta'$ are the angular distances of the planets from a fixed point on the sphere, and $\phi$, $\phi'$ are their position-angles, then
\[
\cos z = \cos\theta \cos\theta' + \sin\theta \sin\theta' \cos(\phi - \phi').
\]
This formula is sometimes useful in the theory of the mutual perturbations of the planets; and it may be employed in the calculation of the disturbing function.

\smallskip

\noindent\textbf{6.}\quad In the theory of the mutual perturbations of the planets, the expression for the angular distance enters into the disturbing function; and it is necessary to develop this expression in a series of cosines of multiples of the mean longitudes. The development is effected by means of the formula
\[
\frac{1}{\Delta} = \frac{1}{\sqrt{r^2 + r'^2 - 2rr' \cos z}},
\]
where $\Delta$ is the distance between the two planets. This expression is expanded in a series of Legendre's coefficients; and the coefficients are then developed in series of cosines of multiples of the mean longitudes.

\smallskip

\noindent\textbf{7.}\quad The calculation of the coefficients of the series for $\dfrac{1}{\Delta}$ is one of the most laborious parts of the theory of planetary perturbations. The coefficients are functions of the ratios of the distances and of the eccentricities and inclinations of the orbits; and they must be calculated to a high degree of accuracy in order to obtain satisfactory results. The labour of the calculation has been greatly diminished by the introduction of mechanical quadratures and other numerical methods.

\newpage
%% ========================================================================
%% PAPER 470: NOTE ON THE ATTRACTION OF ELLIPSOIDS
%% ========================================================================
\section*{470. Note on the Attraction of Ellipsoids.}
\addcontentsline{toc}{section}{470. Note on the Attraction of Ellipsoids (pp.~380--383)}

\noindent\textit{[From the Monthly Notices of the Royal Astronomical Society, vol.~XXV. (1864--1865), pp.~258--264.]}

\medskip

\noindent\textbf{1.}\quad The attraction of an ellipsoid upon an external point is a problem which has occupied the attention of mathematicians since the time of Newton. The problem was first completely solved by Ivory, who showed that the attraction of an ellipsoid upon an external point is the same as the attraction of a certain sphere upon the same point; and this result has been extended by Chasles and others to the case of heterogeneous ellipsoids.

\smallskip

\noindent\textbf{2.}\quad The method of Ivory consists in comparing the attraction of the given ellipsoid with that of a confocal ellipsoid passing through the attracted point. If $a$, $b$, $c$ are the semi-axes of the given ellipsoid, and $a'$, $b'$, $c'$ are those of the confocal ellipsoid, then
\[
\frac{x^2}{a^2} + \frac{y^2}{b^2} + \frac{z^2}{c^2} = 1, \qquad
\frac{x^2}{a'^2} + \frac{y^2}{b'^2} + \frac{z^2}{c'^2} = 1,
\]
are the equations of the two ellipsoids; and the condition that they are confocal gives
\[
a'^2 - a^2 = b'^2 - b^2 = c'^2 - c^2 = \theta,
\]
say, where $\theta$ is determined by the condition that the second ellipsoid passes through the attracted point.

\smallskip

\noindent\textbf{3.}\quad The potentials of the two ellipsoids at the attracted point are
\[
V = \pi\rho abc \int_0^\infty \frac{ds}{\sqrt{(a^2+s)(b^2+s)(c^2+s)}},
\]
\[
V' = \pi\rho a'b'c' \int_0^\infty \frac{ds}{\sqrt{(a'^2+s)(b'^2+s)(c'^2+s)}};
\]
and Ivory's theorem asserts that the attractions of the two ellipsoids at the attracted point are in the ratio of the masses of the ellipsoids. That is, if $X$, $Y$, $Z$ are the components of the attraction of the first ellipsoid, and $X'$, $Y'$, $Z'$ are those of the second, then
\[
\frac{X}{X'} = \frac{Y}{Y'} = \frac{Z}{Z'} = \frac{abc}{a'b'c'}.
\]

\smallskip

\noindent\textbf{4.}\quad This theorem may be proved by direct integration. The potential of a homogeneous ellipsoid at an external point $(x,y,z)$ is
\[
V = \pi\rho abc \int_\lambda^\infty \left(1 - \frac{x^2}{a^2+s} - \frac{y^2}{b^2+s} - \frac{z^2}{c^2+s}\right) \frac{ds}{\sqrt{(a^2+s)(b^2+s)(c^2+s)}},
\]
where $\lambda$ is the greatest root of the equation
\[
\frac{x^2}{a^2+\lambda} + \frac{y^2}{b^2+\lambda} + \frac{z^2}{c^2+\lambda} = 1.
\]

\newpage
\noindent\textbf{5.}\quad The components of the attraction are obtained by differentiating the potential with respect to $x$, $y$, $z$; and we find
\begin{align*}
X &= -\frac{\partial V}{\partial x} = 2\pi\rho abc \cdot x \int_\lambda^\infty \frac{ds}{(a^2+s)\sqrt{(a^2+s)(b^2+s)(c^2+s)}},\\
Y &= -\frac{\partial V}{\partial y} = 2\pi\rho abc \cdot y \int_\lambda^\infty \frac{ds}{(b^2+s)\sqrt{(a^2+s)(b^2+s)(c^2+s)}},\\
Z &= -\frac{\partial V}{\partial z} = 2\pi\rho abc \cdot z \int_\lambda^\infty \frac{ds}{(c^2+s)\sqrt{(a^2+s)(b^2+s)(c^2+s)}}.
\end{align*}
These expressions give the attraction of a homogeneous ellipsoid upon an external point in terms of definite integrals.

\smallskip

\noindent\textbf{6.}\quad For a point on the surface of the ellipsoid, we have $\lambda=0$; and the expressions become
\begin{align*}
X &= 2\pi\rho abc \cdot x \int_0^\infty \frac{ds}{(a^2+s)\sqrt{(a^2+s)(b^2+s)(c^2+s)}},\\
Y &= 2\pi\rho abc \cdot y \int_0^\infty \frac{ds}{(b^2+s)\sqrt{(a^2+s)(b^2+s)(c^2+s)}},\\
Z &= 2\pi\rho abc \cdot z \int_0^\infty \frac{ds}{(c^2+s)\sqrt{(a^2+s)(b^2+s)(c^2+s)}}.
\end{align*}
These integrals may be evaluated in terms of elliptic functions; and the result is the well-known formula for the attraction of an ellipsoid upon a point on its surface.

\smallskip

\noindent\textbf{7.}\quad The attraction of a heterogeneous ellipsoid, in which the density is a function of the distance from the centre, may be obtained by integrating the attraction of a series of concentric homogeneous ellipsoids. If the density is $\rho = f(\sqrt{x^2+y^2+z^2})$, then the potential is
\[
V = \int_0^1 \pi f(ar) \, r^2 dr \int_0^\infty \frac{ds}{\sqrt{(a^2r^2+s)(b^2r^2+s)(c^2r^2+s)}},
\]
and the components of the attraction are obtained by differentiating this expression.

\newpage
%% ========================================================================
%% PAPER 471: ON THE DETERMINATION OF A PLANET'S ORBIT
%% ========================================================================
\section*{471. On the Determination of a Planet's Orbit.}
\addcontentsline{toc}{section}{471. On the Determination of a Planet's Orbit (pp.~384--386)}

\noindent\textit{[From the Monthly Notices of the Royal Astronomical Society, vol.~XXV. (1864--1865), pp.~264--268.]}

\medskip

\noindent\textbf{1.}\quad The determination of the orbit of a planet from three observations is one of the fundamental problems of theoretical astronomy; and it has been treated by many eminent mathematicians, from the time of Gauss to the present day. The problem consists in finding the elements of the orbit, that is, the semi-axis major, the eccentricity, the inclination, the longitude of the node, the longitude of the perigee, and the mean motion, from three observed positions of the planet.

\smallskip

\noindent\textbf{2.}\quad The method of Gauss, which is the most celebrated of all the methods for the solution of this problem, consists in expressing the ratios of the triangles formed by the three observed positions of the planet and the Sun in terms of the elements of the orbit; and then determining the elements by means of these ratios. The method is elegant and complete; and it has been the foundation of all subsequent investigations on the subject.

\smallskip

\noindent\textbf{3.}\quad The fundamental equations of the method of Gauss are the following. Let $r_1$, $r_2$, $r_3$ be the radii vectores of the planet at the three epochs; let $2f_1$, $2f_2$, $2f_3$ be the areas of the triangles formed by the radii vectores; and let $\tau_1$, $\tau_2$, $\tau_3$ be the intervals of time between the observations. Then
\[
\frac{f_1}{\tau_1} = \frac{\sqrt{\mu p}}{2}, \qquad
\frac{f_2}{\tau_2} = \frac{\sqrt{\mu p}}{2}, \qquad
\frac{f_3}{\tau_3} = \frac{\sqrt{\mu p}}{2},
\]
where $\mu$ is the sum of the masses of the Sun and planet, and $p$ is the semi-parameter of the orbit.

\smallskip

\noindent\textbf{4.}\quad The ratios of the areas of the triangles are known from the observations; and the ratios of the intervals of time are also known. The equations therefore give the value of $p$; and this value being substituted in the equations of the orbit, the elements may be determined.

\smallskip

\noindent\textbf{5.}\quad The method of Gauss has been modified and improved by many subsequent investigators; and various methods have been proposed for the direct solution of the equations, without the use of successive approximations. The most important of these methods are those of Encke, Leverrier, and others; and they have been employed with great success in the calculation of the orbits of the planets and comets.

\newpage
%% ========================================================================
%% PAPER 472: NOTE ON LAMBERT'S THEOREM FOR ELLIPTIC MOTION
%% ========================================================================
\section*{472. Note on Lambert's Theorem for Elliptic Motion.}
\addcontentsline{toc}{section}{472. Note on Lambert's Theorem for Elliptic Motion (pp.~387--389)}

\noindent\textit{[From the Monthly Notices of the Royal Astronomical Society, vol.~XXV. (1864--1865), pp.~268--272.]}

\medskip

\noindent\textbf{1.}\quad Lambert's theorem is one of the most remarkable results in the theory of elliptic motion; and it has been the subject of many investigations by eminent mathematicians. The theorem asserts that the time of describing an arc of an ellipse depends only on the sum of the radii vectores of the extremities of the arc, and on the chord of the arc; or, in other words, that elliptic arcs which have the same chord and the same sum of radii vectores are described in the same time.

\smallskip

\noindent\textbf{2.}\quad The analytical expression of Lambert's theorem is as follows. Let $r$, $r'$ be the radii vectores of the extremities of an arc of an ellipse; let $c$ be the chord of the arc; and let $t$ be the time of describing the arc. Then
\[
\sqrt{\mu}\, t = a^{\frac{3}{2}} \left\{ (\epsilon - \sin\epsilon) - (\delta - \sin\delta) \right\},
\]
where $\mu$ is the sum of the masses, $a$ is the semi-axis major, and $\epsilon$, $\delta$ are angles determined by the equations
\[
\sin\frac{\epsilon}{2} = \frac{1}{2}\sqrt{\frac{r+r'+c}{a}}, \qquad
\sin\frac{\delta}{2} = \frac{1}{2}\sqrt{\frac{r+r'-c}{a}}.
\]

\smallskip

\noindent\textbf{3.}\quad The proof of Lambert's theorem may be obtained by a direct integration of the differential equations of elliptic motion. The time of describing an arc is expressed by the integral
\[
t = \int \frac{r^2 \, dv}{\sqrt{\mu p}},
\]
where $v$ is the true anomaly and $p$ is the semi-parameter. By transforming this integral by means of the relations between the true anomaly, the eccentric anomaly, and the mean anomaly, the result may be brought to the form given by Lambert's theorem.

\smallskip

\noindent\textbf{4.}\quad Another proof of Lambert's theorem may be obtained by means of the theory of elliptic functions. The expression for the time involves the elliptic integral of the second kind; and this integral may be transformed by means of the addition-formulae of the elliptic functions. The result of this transformation is Lambert's theorem; and the proof is in many respects simpler and more elegant than the direct proof.

\newpage
\noindent\textbf{5.}\quad Lambert's theorem may be extended to hyperbolic motion. In this case, the time of describing an arc is expressed by means of hyperbolic functions; and the formula is
\[
\sqrt{\mu}\, t = (-a)^{\frac{3}{2}} \left\{ (\sinh\epsilon - \epsilon) - (\sinh\delta - \delta) \right\},
\]
where $\epsilon$, $\delta$ are determined by equations similar to those for the elliptic case, but with hyperbolic functions in place of circular functions.

\smallskip

\noindent\textbf{6.}\quad The theorem may also be extended to parabolic motion. In this case, the semi-axis major is infinite; and the formula becomes
\[
t = \frac{1}{6\sqrt{\mu}} \left\{ (r+r'+c)^{\frac{3}{2}} \pm (r+r'-c)^{\frac{3}{2}} \right\},
\]
which is Barker's theorem. This theorem is a particular case of Lambert's theorem; and it may be obtained from it by making the eccentricity equal to unity.

\smallskip

\noindent\textbf{7.}\quad Lambert's theorem has many important applications in the theory of planetary motion. It is used in the determination of the orbits of planets and comets from observations; and it is also of use in the calculation of the perturbations of the planets. The theorem is one of the most elegant results in the theory of elliptic motion; and it is remarkable for its generality and simplicity.

\newpage
%% ========================================================================
%% PAPER 473: ON THE CONSTRUCTION OF THE UMBRAL OR PENUMBRAL CURVE &C.
%% ========================================================================
\section*{473. On the Construction of the Umbral or Penumbral Curve \&c.}
\addcontentsline{toc}{section}{473. On the Construction of the Umbral or Penumbral Curve (pp.~390--391)}

\noindent\textit{[From the Monthly Notices of the Royal Astronomical Society, vol.~XXV. (1864--1865), pp.~272--274.]}

\medskip

\noindent\textbf{1.}\quad The umbral and penumbral curves are the boundaries of the regions on the surface of the Earth where the Sun is totally or partially eclipsed by the Moon. The construction of these curves is a problem of considerable importance in the theory of eclipses; and it has been treated by many eminent astronomers.

\smallskip

\noindent\textbf{2.}\quad The umbral curve is the locus of points on the surface of the Earth where the Sun is totally eclipsed; and the penumbral curve is the locus of points where the Sun is partially eclipsed. These curves are determined by the condition that the apparent distances of the centres of the Sun and Moon, as seen from a point on the Earth's surface, are equal to the sum or difference of their apparent semidiameters.

\smallskip

\noindent\textbf{3.}\quad The condition for the umbral curve is
\[
\Delta = s - s',
\]
where $\Delta$ is the apparent distance of the centres of the Sun and Moon, $s$ is the apparent semidiameter of the Sun, and $s'$ is the apparent semidiameter of the Moon. The condition for the penumbral curve is
\[
\Delta = s + s'.
\]

\smallskip

\noindent\textbf{4.}\quad The construction of these curves on the Earth's surface may be effected by the following method. Let $S$ and $M$ be the centres of the Sun and Moon; let $O$ be the centre of the Earth; and let $P$ be a point on the surface of the Earth. The apparent distance of the centres of the Sun and Moon, as seen from $P$, is determined by the spherical triangle $SPM$; and the condition for the umbral or penumbral curve gives an equation between the coordinates of $P$.

\smallskip

\noindent\textbf{5.}\quad The equation of the umbral or penumbral curve may be written in the form
\[
\cos SP \cos MP + \sin SP \sin MP \cos SPM = \cos(s \mp s').
\]
This equation determines the locus of the point $P$ on the surface of the Earth; and it may be solved by expressing the coordinates of $P$ in terms of the latitude and longitude.

\newpage
%% ========================================================================
%% PAPER 474: ON THE GEOMETRICAL THEORY OF SOLAR ECLIPSES
%% ========================================================================
\section*{474. On the Geometrical Theory of Solar Eclipses.}
\addcontentsline{toc}{section}{474. On the Geometrical Theory of Solar Eclipses (pp.~392--396)}

\noindent\textit{[From the Monthly Notices of the Royal Astronomical Society, vol.~XXV. (1864--1865), pp.~274--282.]}

\medskip

\noindent The geometrical theory of solar eclipses is founded upon the consideration of the cones of shadow cast by the Moon. These cones are determined by the apparent semidiameters of the Sun and Moon, and by their relative positions; and their intersections with the surface of the Earth give the regions of total and partial eclipse.

\smallskip

\noindent The cone of total shadow, or the umbra, is the cone whose vertex is between the Moon and the Earth, and whose generators touch the Sun and Moon externally. The cone of partial shadow, or the penumbra, is the cone whose vertex is on the side of the Moon remote from the Earth, and whose generators touch the Sun and Moon internally on one side and externally on the other.

\smallskip

\noindent The intersections of these cones with the surface of the Earth are curves of a complicated character; and their determination requires the solution of a system of equations involving the coordinates of the Sun, Moon, and Earth. The following investigation is founded upon a method which was first employed by Bessel, and which has been adopted by most subsequent writers on the subject.

\smallskip

\noindent Let $\xi$, $\eta$, $\zeta$ be the coordinates of the Moon relative to the centre of the Earth; let $R$ be the distance of the Moon from the Earth; and let $s$, $s'$ be the apparent semidiameters of the Sun and Moon. Then the equations of the cones of shadow are
\begin{align*}
\xi^2 + \eta^2 &= (\zeta \tan s' - R \sin s)^2 \quad \text{(umbra)},\\
\xi^2 + \eta^2 &= (\zeta \tan s' + R \sin s)^2 \quad \text{(penumbra)}.
\end{align*}

\smallskip

\noindent The intersections of these cones with the surface of the Earth are obtained by substituting for $\xi$, $\eta$, $\zeta$ their values in terms of the latitude and longitude of the point on the Earth's surface. The resulting equations are of a complicated character; but they may be solved by numerical methods.

\newpage
\noindent\textbf{1.}\quad The fundamental equations of the geometrical theory of eclipses are the following. Let $x$, $y$, $z$ be the coordinates of the Moon relative to the centre of the Earth; let $l$ be the ratio of the apparent semidiameter of the Moon to the apparent semidiameter of the Sun; and let $\sin f$, $\sin f'$ be the sines of the angles which the generators of the umbral and penumbral cones make with the line joining the centres of the Sun and Moon. Then the equations of the umbral and penumbral surfaces are
\begin{align*}
(x-x_0)^2 + (y-y_0)^2 &= (z-z_0)^2 \tan^2 f \quad \text{(umbra)},\\
(x-x_0)^2 + (y-y_0)^2 &= (z-z_0)^2 \tan^2 f' \quad \text{(penumbra)},
\end{align*}
where $x_0$, $y_0$, $z_0$ are the coordinates of the vertex of the cone.

\smallskip

\noindent\textbf{2.}\quad The coordinates of the vertex of the umbral cone are determined by the condition that the distances of the vertex from the centres of the Sun and Moon are in the ratio of their apparent semidiameters. If $D$ is the distance between the centres of the Sun and Moon, then
\[
\frac{z_0}{R'} = \frac{s'}{s-s'}, \qquad z_0 = \frac{D s'}{s-s'},
\]
where $R'$ is the distance of the Sun from the Earth. The coordinates $x_0$, $y_0$ are equal to the coordinates of the Moon multiplied by the ratio $\dfrac{z_0}{R}$.

\smallskip

\noindent\textbf{3.}\quad The intersections of the umbral and penumbral surfaces with the surface of the Earth are the curves of total and partial eclipse. These curves may be constructed by the following method. For a given time, calculate the coordinates of the Moon and the values of $s$, $s'$; then calculate the coordinates of the vertex of the cone; and finally, determine the intersection of the cone with the sphere of the Earth.

\smallskip

\noindent\textbf{4.}\quad The equation of the intersection of the cone with the sphere of the Earth may be written in the form
\[
\xi^2 + \eta^2 + \zeta^2 = R_\oplus^2, \qquad
\xi^2 + \eta^2 = (\zeta - z_0)^2 \tan^2 f,
\]
where $R_\oplus$ is the radius of the Earth, and $\xi$, $\eta$, $\zeta$ are the coordinates of a point on the curve. Eliminating $\xi^2 + \eta^2$, we obtain
\[
R_\oplus^2 - \zeta^2 = (\zeta - z_0)^2 \tan^2 f,
\]
which is a quadratic equation for $\zeta$.

\newpage
\noindent\textbf{5.}\quad The solution of this equation gives two values for $\zeta$; and each value corresponds to a point on the curve. The values of $\xi$ and $\eta$ are then determined by the equation of the cone; and the latitude and longitude of the point are obtained by transforming the rectangular coordinates to geographical coordinates.

\smallskip

\noindent\textbf{6.}\quad The curve of total eclipse is thus determined for each instant of time; and the path of the eclipse on the Earth's surface is the locus of these curves. The limits of the path are determined by the condition that the quadratic equation for $\zeta$ has equal roots; and this condition gives the equations of the limiting curves.

\smallskip

\noindent\textbf{7.}\quad The duration of the total eclipse at any point on the path is determined by the time during which the point remains within the umbral cone. This time may be calculated from the motion of the Moon and the rotation of the Earth; and it is a maximum at the centre of the path, and decreases to zero at the limits.

\smallskip

\noindent\textbf{8.}\quad The foregoing theory may be extended to the case of annular eclipses. In this case, the vertex of the umbral cone is between the Earth and the Moon; and the eclipse is annular when the apparent diameter of the Moon is less than that of the Sun. The condition for an annular eclipse is that the umbral cone intersects the surface of the Earth beyond the vertex; and the theory of annular eclipses is in all respects similar to that of total eclipses.

\smallskip

\noindent\textbf{9.}\quad The penumbral cone gives rise to partial eclipses; and the region of partial eclipse is much more extensive than the region of total eclipse. The limits of the partial eclipse are determined by the intersection of the penumbral cone with the surface of the Earth; and they may be calculated by the same method as the limits of the total eclipse.

\newpage
%% ========================================================================
%% PAPER 475: ON A PROPERTY OF THE STEREOGRAPHIC PROJECTION
%% ========================================================================
\section*{475. On a Property of the Stereographic Projection.}
\addcontentsline{toc}{section}{475. On a Property of the Stereographic Projection (pp.~397--399)}

\noindent\textit{[From the Monthly Notices of the Royal Astronomical Society, vol.~XXV. (1864--1865), pp.~282--286.]}

\medskip

\noindent The stereographic projection is one of the most important of the methods of representing the surface of a sphere upon a plane. It possesses many remarkable properties; among which may be mentioned the property that circles on the sphere are projected into circles on the plane, and the property that the projection is conformal, that is, that angles are preserved unaltered.

\smallskip

\noindent The property which I now proceed to investigate relates to the representation of great circles upon the stereographic projection. A great circle on the sphere is projected into a circle on the plane which passes through the projections of the poles of the great circle; and the centre of this circle is the projection of the pole of the great circle. The radius of the projected circle depends upon the distance of the great circle from the pole of projection; and it may be expressed in terms of this distance by a simple formula.

\smallskip

\noindent Let $P$ be the pole of projection, and let $C$ be the centre of the sphere. Let $Q$ be the pole of a great circle on the sphere, and let $R$ be the radius of the sphere. Then the distance of $Q$ from $P$ is $PQ = R \cdot \theta$, where $\theta$ is the angular distance; and the projection of the great circle is a circle whose centre is the projection of $Q$, and whose radius is $R \tan\frac{\theta}{2}$.

\smallskip

\noindent This property may be proved as follows. Let $AB$ be the great circle, and let $Q$ be its pole. Draw the plane through $P$, $C$, and $Q$; and let this plane intersect the great circle $AB$ in the points $A$ and $B$. Then $PA$ and $PB$ are the generators of the cone which projects the great circle; and the intersection of this cone with the plane of projection is the projected circle.

\newpage
\noindent\textbf{1.}\quad The stereographic projection of a sphere upon a plane is effected by drawing lines from a fixed point on the surface of the sphere to the points of the sphere, and taking the intersections of these lines with a fixed plane. The fixed point is called the pole of projection; and the fixed plane is called the plane of projection.

\smallskip

\noindent\textbf{2.}\quad The most important case is that in which the plane of projection is perpendicular to the diameter through the pole of projection. In this case, the projection is called the orthogonal stereographic projection; and it possesses the property that circles on the sphere are projected into circles on the plane.

\smallskip

\noindent\textbf{3.}\quad The formula for the radius of the projected circle of a great circle may be obtained as follows. Let $\theta$ be the angular distance of the great circle from the pole of projection; then the radius of the projected circle is
\[
r = R \tan\frac{\theta}{2},
\]
where $R$ is the radius of the sphere. This formula may be proved by considering the right-angled triangle formed by the pole of projection, the pole of the great circle, and a point on the great circle.

\smallskip

\noindent\textbf{4.}\quad The centres of the projected circles of all great circles lie on the plane of projection; and they are the projections of the poles of the great circles. The locus of the centres of the projected circles is therefore the whole plane of projection; and every point of this plane is the centre of the projection of some great circle.

\smallskip

\noindent\textbf{5.}\quad The property that the stereographic projection is conformal may be proved by showing that the scale of the projection is the same in all directions at any point. The scale of the projection is given by the formula
\[
\frac{ds'}{ds} = \frac{1}{\cos^2\frac{\theta}{2}},
\]
where $ds$ is an element of arc on the sphere, $ds'$ is the corresponding element on the plane, and $\theta$ is the angular distance from the pole of projection. Since this expression depends only on $\theta$, and not on the direction of the element, the projection is conformal.

\smallskip

\noindent\textbf{6.}\quad The stereographic projection is extensively used in astronomy for the construction of maps of the celestial sphere; and it is also used in navigation for the construction of Mercator's chart. The property that circles are projected into circles renders the projection peculiarly adapted for the representation of the celestial sphere; and the conformal property renders it useful for the measurement of angles.

\newpage
%% ========================================================================
%% PAPER 476 (start): ON THE DETERMINATION OF THE ORBIT OF A PLANET
%% ========================================================================
\section*{476. On the Determination of the Orbit of a Planet from Three Observations.}
\addcontentsline{toc}{section}{476. On the Determination of the Orbit of a Planet from Three Observations (p.~400, start)}

\noindent\textit{[From the Memoirs of the Royal Astronomical Society, vol.~XXXVIII. (1870), pp.~17--111. Read December 10, 1869.]}

\medskip

\noindent I propose to consider from a geometrical point of view the problem of the determination of the orbit of a planet from three observations. The problem has been treated by many eminent mathematicians; and the methods which have been given are in general analytical. The geometrical method has the advantage of presenting the results in a more intuitive form; and it affords a useful check upon the analytical investigations.

\smallskip

\noindent The three observations give the directions of the planet, as seen from the Earth, at three different epochs; and the problem is to determine the orbit which passes through the three points thus determined. The orbit is an ellipse with the Sun in one focus; and the elements of the orbit are the semi-axis major, the eccentricity, the inclination, the longitude of the node, the longitude of the perigee, and the mean motion.

\smallskip

\noindent The geometrical construction is founded upon the following principles. The three lines of sight from the Earth to the planet determine three points on the celestial sphere; and the problem is to find an ellipse, with the Sun in one focus, which passes through the three points determined by these lines. The construction depends upon the properties of the cone which has the Sun for its vertex, and which passes through the three positions of the planet.

\smallskip

\noindent Let $S$ be the Sun, $E_1$, $E_2$, $E_3$ the three positions of the Earth, and $P_1$, $P_2$, $P_3$ the three positions of the planet. The lines $E_1P_1$, $E_2P_2$, $E_3P_3$ are the lines of sight; and the problem is to determine the ellipse with focus $S$ which passes through $P_1$, $P_2$, $P_3$. The cone with vertex $S$ and generators $SP_1$, $SP_2$, $SP_3$ intersects the celestial sphere in a conic; and the problem is reduced to the determination of this conic.

\smallskip

\noindent The investigation will be continued in the subsequent pages of this memoir; and the complete solution of the problem will be given, with numerical examples and applications to the orbits of the planets.

\vfill
\begin{center}
\rule{0.5\textwidth}{0.4pt}
\end{center}

\end{document}
